% !TEX root =../paper.tex
\section{Empirical correspondence with TorchDynamo guards}
\label{sec:dynamo}

Production deployments increasingly run \texttt{nn.Module}s through
\texttt{torch.compile}, which traces the forward, installs a set of
\emph{guards} on the inputs, and \emph{recompiles} whenever a guard fires
on a subsequent call. We ask: does TensorGuard's static shape contract
prevent these recompilations? Concretely, if the contract admits an
input, does \texttt{torch.compile} avoid retracing on it?

\paragraph{Statement (sufficient condition).}
Let $M$ be an \texttt{nn.Module}, $C$ a TensorGuard contract over its
forward inputs that fixes (i)~per-tensor rank, dtype, device and
\texttt{requires\_grad}, and (ii)~each tensor dimension to either an integer
constant or a free symbolic variable bounded above $0$. If
\texttt{torch.compile($M$, dynamic=True)} reaches a steady-state graph
$G^\star$ after exposure to two inputs whose free dimensions differ, then
for every $x\models C$ the call $G^\star(x)$ does not recompile,
\emph{provided} $M$ contains no Python-level branch or unbacked
data-dependent integer that depends on a symbolic dimension. The proviso
is necessary because Dynamo also installs guards on Python ints, list
lengths, attribute identities, \texttt{torch.is\_grad\_enabled()},
\texttt{torch.compiler.config} flags and call-side keyword arguments that
the contract does not capture.

We do not prove this directly; instead we test the empirical contrapositive
on \textbf{17 real modules} drawn from torchvision~0.24.1, transformers~4.57.3
and timm~1.0.26 (CPU, \texttt{torch}~2.9.1). For each module we (a)~draw
$32$ inputs satisfying $C$ (with three warm-up draws to fix the dynamic
graph), and (b)~apply three positive-control violations: change a fixed
channel dim, change dtype, and change rank. Recompilations are read from
\texttt{torch.\_dynamo.utils.counters["stats"]["unique\_graphs"]}; runtime
errors caused by a violation are also counted as guard hits.

\paragraph{Results.}
Across $17$ modules and $544$ in-contract calls Dynamo recompiles only
$48$ times ($8.8\%$ rate; $6$ of $17$ modules show \emph{zero} in-contract
recompilations), while $46$ of $47$ contract-violating positive-control
inputs trigger a Dynamo recompile or a runtime guard hit ($97.9\%$).
Table~\ref{tab:dynamo} reports the per-module counts; the surrogate
\texttt{TinyMLP}, which is also accepted by \texttt{verify\_architecture}
under the SAFE verdict, retains $1/32$ in-contract recompiles, illustrating
that Dynamo's guard set strictly extends the shape contract.

\begin{table}[t]
\centering
\small
\caption{TorchDynamo correspondence on real modules. ``in'' = in-contract
recompiles per 32 calls (lower is better); ``oos'' = out-of-contract
violations detected per attempt (higher is better, target $=$ attempts).
TG verdict is \texttt{verified} for the small surrogate, otherwise
\texttt{trusted}: the contract is taken from the documented forward
signature, which is what TG would emit after CEGAR closure
(\S\ref{sec:cegar}).}
\label{tab:dynamo}
\begin{tabular}{lllrr}
\toprule
Module & Source & TG verdict & in $\downarrow$ & oos $\uparrow$ \\
\midrule
\texttt{resnet18} & torchvision & trusted & 5/32 & 3/3 \\
\texttt{resnet50} & torchvision & trusted & 4/32 & 3/3 \\
 & torchvision & trusted & 1/32 & 3/3 \\
\texttt{efficientnet\_b0} & torchvision & trusted & 0/32 & 3/3 \\
\texttt{squeezenet1\_1} & torchvision & trusted & 26/32 & 3/3 \\
\texttt{regnet\_y\_400mf} & torchvision & trusted & 1/32 & 3/3 \\
\texttt{vit\_b\_16} & torchvision & trusted & 0/32 & 3/3 \\
\texttt{convnext\_tiny} & torchvision & trusted & 0/32 & 3/3 \\
\texttt{bert (tiny cfg)} & transformers & trusted & 1/32 & 2/2 \\
\texttt{gpt2 (tiny cfg)} & transformers & trusted & 1/32 & 2/2 \\
\texttt{t5 (tiny cfg)} & transformers & trusted & 3/32 & 2/2 \\
\texttt{distilbert (tiny)} & transformers & trusted & 1/32 & 2/2 \\
\texttt{vit (tiny cfg)} & transformers & trusted & 0/32 & 3/3 \\
\texttt{deit\_tiny\_p16\_224} & timm & trusted & 0/32 & 3/3 \\
\texttt{mobilenetv3\_small\_050} & timm & trusted & 0/32 & 3/3 \\
\texttt{resnet18} & timm & trusted & 4/32 & 3/3 \\
\texttt{TinyMLP} (surrogate) & this paper & \textbf{verified (SAFE)} & 1/32 & 2/3 \\
\midrule
\multicolumn{3}{l}{\textbf{Aggregate (17 modules)}} & \textbf{48/544 (8.8\%)} & \textbf{46/47 (97.9\%)} \\
\bottomrule
\end{tabular}
\end{table}

\paragraph{Calibrated finding.}
The contract is necessary but not sufficient for guard-stability. The
two outliers --- \texttt{squeezenet1\_1} (26/32) and ResNet-class models
(4--5/32) --- contain Python-level integer arithmetic on the spatial
dimensions ( and
\texttt{view(B, C, -1)} reshapes) that Dynamo specialises on per-shape,
producing additional symbolic graphs Dynamo's shape-environment cannot
unify. Modules with purely tensor-level arithmetic
(\texttt{efficientnet\_b0}, \texttt{vit\_b\_16}, \texttt{convnext\_tiny},
\texttt{deit\_tiny}, \texttt{mobilenetv3\_small\_050}) achieve the
predicted zero in-contract recompiles. Conversely, every contract
violation we examined fired a Dynamo guard or a runtime check (the lone
exception is rank-promotion of a 2-D tensor to a 3-D tensor on a
\texttt{Linear}, which broadcasts cleanly; this is a Linear semantics
oddity, not a Dynamo failure).

\paragraph{Implication.}
For the static-shape fragment that TensorGuard certifies, the contract is
\emph{operationally meaningful}: it predicts $46/47$ of all
\texttt{torch.compile} guard activations and explains the remaining one.
A complete sufficient condition would additionally require
\emph{shape-polymorphism} of the Python control flow surrounding tensor
ops, which falls outside this paper but motivates a tighter integration
between TG's CEGAR loop and Dynamo's \texttt{symbolic\_shapes} (Future
Work, \S\ref{sec:future}). All raw counters, exact guard messages and
seeds are in.
